\subsection{The residue map}\label{rl-000V}\label{sec:the-residue-map}

\red{Nothing here requires localizing at \(Q_T\) in the domain}

We now define a version of the residue map which will serve as the bridge between invariants of \(\mathcal X\) and invariants of \(\mathcal X^T\). We'll prove it for an open and closed substack \(\mathcal Z\subseteq \mathcal X^T\), open so that we may interchange \(\tau_{\geq -1}L_\eta|_{\mathcal Z}\) with \(\tau_{\geq -1}L_{\eta|_{\mathcal Z}}\) and closed so that the inclusion \(\mathcal Z\hookrightarrow \mathcal X^T\) is a closed immersion. We abuse notation and write \(\eta\) for both \(\eta:\mathcal X^T\to \mathcal X\) and \(\eta|_{\mathcal Z}:\mathcal Z\to \mathcal X\).

Taking inspiration from the Edidin-Graham case in Proposition~\ref{prop:edidin-graham-residue-identity}, the map \(\Res_{\mathcal Z/\mathcal X}:A_*^T(\mathcal X)\to A_*^T(\mathcal Z)_{\loc}\) we wish to write down should be
\begin{align}\label{eqn:preliminary-Res-defn}
  \Res_{\mathcal Z/\mathcal X}(\alpha) = e_T^{-1}(\mathfrak E)\cap \eta^!_{\mathfrak E}(\alpha)
\end{align}
where \(\eta^!_{\mathfrak E}\) is virtual pullback. Virtual pullback requires the specification of additional data, namely, \((j,\mathfrak E)\) where \(\mathfrak E\) is a vector bundle stack over \(\mathcal Z\) and \(j:\mathfrak C_{\mathcal Z/\mathcal X}\hookrightarrow \mathfrak E\) is a closed embedding of the intrinsic normal cone of \(\eta:\mathcal Z\to \mathcal X\). For this to make sense equivariantly we additionally need that \(\mathfrak E\) is \(T\)-equivariant, and to ensure \(e_T(\mathfrak E)\) is a unit we require that \(\mathfrak E\) is moving. We'll refer to $(j,\mathfrak E)$ as \emph{embedding data}.

Given a choice of embedding data \((j, \mathfrak E)\) for the cone stack \(\mathfrak C_\eta\), the equivariant virtual pullback \(\eta^!\) \cite{manolache_VirtualPullbacks2012} is the composite
\begin{align*}
  \eta^!:A^T_*(\mathcal X) \xrightarrow{\sigma} A^T_*(\mathfrak C_\eta) \xrightarrow{j_*} A^T_*(\mathfrak E) \xrightarrow{0^!_{\mathfrak E}} A^T_*(\mathcal Z)
\end{align*}
where \(\sigma\) is the equivariant specialization map on \(\mathcal Z \to \mathcal X\) \(j_*\) is equivariant projective pushforward \(0^!_{\mathfrak E}\) is the Gysin map for vector bundle stacks. We make a comment about this last map.

\noindent{Convention (Gysin maps for vector bundle stacks)} If \(\pi_{\mathfrak E}:\mathfrak E\to \mathcal W\) is a vector bundle stack over a stack \(\mathcal W\) which admits a stratification by global quotients, then \(\pi^*_{\mathfrak E}\) is an isomorphism by \cite[Prop. 5.3.2]{kresch_CycleGroupsArtin1999}, and we set \(0^!_{\mathfrak E} := (\pi^*_{\mathfrak E})^!\). Our standing hypotheses include the stratification hypothesis for this reason.


Under our running assumptions, we claim the map~\eqref{eqn:preliminary-Res-defn}~exists and is independent of the choice of embedding data. Existence is the content of Lemma~\ref{lem:existence-of-pot-for-fixed-locus-map}~and independence the content of Lemma~\ref{lem:independence-of-Res-map-from-choice}. Once those are established we define the residue map exactly as above in Definition~\ref{defn:residue-map}.

\begin{remark}\label{rl-000W}\label{rmk:embedding-data-is-POT}
  Before proceeding to the technical lemmas, we remind the reader that specifying (moving \(T\)-equivariant) embedding data is equivalent to specifying a (moving \(T\)-equivariant) perfect obstruction theory \cite[Prop. 3.11]{manolache_VirtualPullbacks2012}, and we will use them interchangeably. If given such a perfect obstruction theory, then the corresponding pair \((j,\mathfrak E)\) is given by \(\mathfrak E = h^1/h^0\big((E^\bullet)^\vee\big)\) and the composition
  \begin{align*}
    j: \mathfrak C_{\mathcal Z/\mathcal X} \hookrightarrow \mathfrak N_{\mathcal Z/\mathcal X} = h^1/h^0(\tau_{\geq -1}L_{\eta|_{\mathcal Z}}^\vee) \overset{~\phi^\vee}{\hookrightarrow}\mathfrak E.
  \end{align*}
\end{remark}


\begin{lemma}[(Existence of embedding data.)]\label{rl-000X}\label{lem:existence-of-pot-for-fixed-locus-map}
  Let \(\mathcal X\) be an algebraic stack finite type over \(k\) with a \(T\)-action and suppose that \(\mathcal X^T\) is algebraic. Let \(\mathcal Z\subseteq \mathcal X^T\) be an open and closed substack which is finite type over \(k\) and has the resolution property (every coherent sheaf on \(\mathcal Z\) is a quotient of finite-rank locally free sheaves). Then there exists a \(T\)-equivariant perfect obstruction theory
  \begin{align*}
    \phi:E^\bullet = [E^{-1} \to E^0] \to \tau = \tau_{\geq -1} L_\eta|_{\mathcal Z}
  \end{align*}
  with \(E^{-1}\), \(E^0\) finite-rank moving \(T\)-equivariant locally free sheaves on \(\mathcal Z\). Consequently
  \begin{align*}
    \mathfrak E := h^1/h^0\big((E^\bullet)^\vee\big) = [E^1/E^0]
  \end{align*}
  is a vector bundle stack over \(\mathcal Z\), globally presented by finite-rank moving locally free sheaves, and there is a \(T\)-equivariant closed embedding of cone stacks
  \begin{align*}
    j:\mathfrak C_{\mathcal Z/\mathcal X} \hookrightarrow \mathfrak E.
  \end{align*}
  In particular, if every connected component of \(\mathcal X^T\) is finite type over \(k\) and has the resolution property, then the conclusion holds over all of \(\mathcal X^T\) componentwise. 
\end{lemma}
\begin{proof} 
  Since \(\eta\) is representable by algebraic spaces it is DM-type (\ref{lem:torus-actions-are-nice}), so the intrinsic normal sheaf and cone \(\mathfrak C_{\mathcal Z/\mathcal X}\subseteq \mathfrak N_{\mathcal Z/\mathcal X}\) are defined and the cone is a subcone of the sheaf \cite[Prop. 2.33]{manolache_VirtualPullbacks2012}. Moreover \(h^1(L_\eta) = 0\) (by \cite[Remark 2.3]{manolache_VirtualPullbacks2012} or \cite[Proposition 2.24]{aranha-pstragowski_IntrinsicNormalCone2024} or \cite[Prop. 2.24]{aranha-pstragowski_IntrinsicNormalCone2024}) so \(\tau\) has cohomology only in degrees \(-1\) and \(0\), and
  \begin{align*}
    \mathfrak N_{\mathcal Z/\mathcal X} = h^1/h^0(\tau^\vee).
  \end{align*}

  \bigskip
  
  \noindent\emph{Step 1: coherent representative for \(\tau\).}~We now argue that \(\tau\) is quasi-isomorphic to a moving complex \([A^{-1}\to A^0]\) with \(A^{-1}\) and \(A^0\) coherent. A quasi-coherent \(T\)-equivariant sheaf on \(\mathcal Z\) is the same as a quasi-coherent sheaf on \(\mathcal Z\times BT\), so all sheaves in this step are taken to live on the Noetherian algebraic stack \(\mathcal Z\times BT\).
  
  Note first that \(h^{-1}\) and \(h^0\) are coherent. Indeed, coherence can be checked locally, and since Olsson's cotangent complex is compatible with restriction along smooth charts, we can reduce this claim to a morphism of algebraic spaces locally of finite type over \(BT\). Here, \(\tau_{\geq -1}L\) is computed by the naive cotangent complex \cite[08RB]{stacks-project}, whose \(h^0\) and \(h^{-1}\) are coherent \cite[0D0U]{stacks-project}.

  Next, because \(\mathcal Z\times BT\) is Noetherian and \(\tau\) is bounded with quasi-coherent cohomology, be can find a representative \(B^\bullet = [B^{-1}\xrightarrow{d} B^0]\) for \(\tau\). Choose \(A^0 \subseteq B^0\) to be a coherent subsheaf surjecting onto \(h^0(\tau)\) under \(B^0\to h^0(\tau)\), this must exist by the coherence of \(h^0(\tau)\). The submodule \(d^{-1}(A^0)\subseteq B^{-1}\) is only quasi-coherent a priori, but the restriction of \(d\) to \(d^{-1}(A^0)\) gives us the exact sequence
  \begin{align*}
    0 \to \ker(d) \to d^{-1}(A^0) \to \img(d)\cap A^0 \to 0,
  \end{align*}
  and since the left and right terms are coherent by the coherence of \(h^0\) and \(h^{-1}\), the middle term is as well. Thus \(A^{-1}:= d^{-1}(A^0)\) is coherent and the inclusion \([A^{-1}\to A^0]\hookrightarrow [B^{-1}\to B^0]\) is a quasi-isomorphism. By Lemma~\ref{lem:two-term-truncation-of-fixed-locus-cotangent-complex-no-fixed-part}, \(L_{\mathcal Z/\mathcal X}\) is entirely moving, and hence we can choose \(A^\bullet\) to have no fixed part.

  \bigskip
  
  \noindent\emph{Step 2: resolve.}~We now resolve \(A^\bullet\) with finite-rank locally free sheaves, and because we assumed the resolution property on \(\mathcal Z\), we once again take our sheaves to live on \(\mathcal Z\) rather than on \(\mathcal Z\times BT\). To ensure the resolution of \(A^\bullet\) is also moving, do this weight by weight: if \(F\) is any coherent moving sheaf on \(\mathcal Z\), then it can be decomposed as \(F = F_{\chi_1}\oplus \cdots \oplus F_{\chi_i}\) for finitely many non-zero weights \(\chi_i\). Then let \(W_i\twoheadrightarrow F_{\chi_i}\) be a resolution of \(F_{\chi_i}\), and let \(T\) act on \(W_i\) with pure weight \(\chi_i\). Then \(\bigoplus_i W_i\to F\) is equivariant, finite-rank and moving.

  We can therefore choose a moving finite-rank locally free surjection \(q_0:E^0\twoheadrightarrow A^0\). Then take
  \begin{align*}
    P := E^0\times_{A^0}A^{-1} = \{(e,a)~\mid~q_0(e) = d(a)\} \subseteq E^0\oplus A^{-1},
  \end{align*}
  which is also moving and coherent, and choose \(q_{-1}:E^{-1}\twoheadrightarrow P\) to be another moving-finite rank locally free surjection. This gives us a complex \(E^\bullet = [E^{-1}\xrightarrow{d_E} E^0]\) and a map \(\phi:E^\bullet \to A^\bullet\) so that \(h^0(\phi)\) is an isomorphism and \(h^{-1}(\phi)\) is a surjection, hence a $T$-equivariant moving perfect obstruction theory for \(\eta\). By Remark~\ref{rmk:embedding-data-is-POT}, this is the same as moving \(T\)-equivariant embedding data for \(\eta\).

\end{proof}

We now state the independence lemma.
\begin{lemma}[Independence of the residue map]\label{rl-000Y}\label{lem:independence-of-Res-map-from-choice}
  Let \(\mathcal X\), \(\mathcal X^T\) and \(\mathcal Z\subseteq\mathcal X^T\) be as in Lemma~\ref{lem:existence-of-pot-for-fixed-locus-map}. Given any two pairs \((j_1, \mathfrak E_1)\) and \((j_2,\mathfrak E_2)\) where \(j_i:\mathfrak C_{\mathcal Z/\mathcal X}\hookrightarrow \mathfrak E_i\) is a \(T\)-equivariant closed embedding and \(\mathfrak E_i\) is a moving \(T\)-equivariant finite rank vector bundle stack for each \(i=1,2\), then for any class \(\alpha\in A_*^T(\mathcal X)\),
  \begin{align*}
    e_T^{-1}(\mathfrak E_1)\cap \eta_1^!(\alpha) = e_T^{-1}(\mathfrak E_2)\cap \eta_2^!(\alpha)
  \end{align*}
  in \(A^T_*(\mathcal X^T)_{\loc}\). Here, \(\eta^!_i\) is the virtual pullback defined with respect to \((j_i,\mathfrak E_i)\).
\end{lemma}


\begin{remark}\label{rl-000Z}\label{rmk:two-modifications-after-the-residue-independence-lemma-statement}
  Before proceeding, two remarks are in order. First, the same specialization map \(\sigma:A^T_*(\mathcal X)\to A^T_*(\mathfrak C)\) occurs in both \(\eta^!_1\) and \(\eta^!_2\), hence it suffices to prove that for every \(\beta\in A^T_*(\mathfrak C_{\mathcal Z/\mathcal X})\),
  \begin{align}\label{eqn:independence-lemma-simplification}
    e_T(\mathfrak E_1)^{-1}\cap 0^!_{\mathfrak E_1}\big((j_1)_*\beta\big) = e_T(\mathfrak E_2)^{-1}\cap 0^!_{\mathfrak E_2}\big((j_2)_*\beta\big).
  \end{align}
  This removes the stack \(\mathcal X\) and the map \(\sigma\) from the discussion entirely.

  Second, once the specialization map \(\sigma\) is removed, the cone structure of \(\mathfrak C\) is no longer relevant, and we can simply take \(\mathfrak C\) to be any algebraic stack \(T\)-equivariant over \(\mathcal Z\) admitting a \(T\)-equivariant closed embedding into \(\mathfrak E\). This is not a superfluous point; the proof of Lemma~\ref{lem:independence-of-Res-map-from-choice} requires this generality.
\end{remark}

For these reasons, we make the following definition.
\begin{definition}\label{rl-0010}
  If \(\mathfrak C\) is an algebraic stack \(T\)-equivariant over \(\mathcal Z\) admitting a \(T\)-equivariant closed embedding \(j:\mathfrak C\hookrightarrow \mathfrak E\) into a finite-rank moving vector bundle stack over \(\mathcal Z\), then for any class \(\beta \in A^T_*(\mathfrak C)\) we define
  \begin{align*}
    S_{(j,\mathfrak E)}(\beta) := e_T(\mathfrak E)^{-1}\cap 0^!_{\mathfrak E}(j_*\beta)\in A^T_*(\mathcal Z)_{\loc}.
  \end{align*}
\end{definition}


We first prove Lemma~\ref{lem:independence-of-Res-map-from-choice}~in the non-stacky case.

\begin{lemma}[Residue independence for honest vector bundles]\label{rl-0011}\label{lem:independence-for-honest-vector-bundles}
  Let \(C\) be a \(T\)-stack over \(\mathcal Z\) admitting \(T\)-equivariant closed embeddings \(j_1:C\hookrightarrow E_1\) and \(j_2:C\hookrightarrow E_2\) into two different \(T\)-equivariant finite rank moving vector bundles over \(\mathcal Z\). Then for all \(\beta \in A^T_*(C)\),
  \[S_{(j_1,E_1)}(\beta) = S_{(j_2,E_2)}(\beta).\]
  We can therefore drop the reference to the embedding data and write \(S_C(\beta)\) without ambiguity.
\end{lemma}
\begin{proof}
  The strategy is to embed \(C\) into \(E_1\oplus E_2\) in two ways: first via
  \begin{align*}
    i_0:C\overset{j_1}{\hookrightarrow} E_1 = E_1\oplus 0 \overset{\iota}{\hookrightarrow} E_1 \oplus E_2,
  \end{align*}
  and then via the diagonal embedding
  \begin{align*}
    i_1:C\longhookrightarrow E_1\oplus E_2, \quad c\mapsto (j_1(c), j_2(c)).
  \end{align*}
  First define the interpolation map
  \begin{align*}
    \varphi:C\times \mathbb A^1 \to E_1\oplus E_2 \times \mathbb A^1; \quad \varphi(c, t) = (j_1(c), t\cdot j_2(c), t),
  \end{align*}
  noting that the fibers of \(\varphi\) at \(C\times\{0\}\) and \(C\times \{1\}\) recover \(i_0\) and \(i_1\) respectively. This map is \(T\)-equivariant if we let \(T\) act trivially on \(\mathbb A^1\).

  Now let \(\pr_C:C\times \mathbb A^1 \to C\) and \(\pr_{\mathcal Z}:\mathcal Z\times \mathbb A^1\to \mathbb Z\) be the projection maps and for \(\beta \in A^T_*(C)\) set
  \begin{align*}
    \Theta(\beta) :&= 0^!_{(E_1\oplus E_2) \times \mathbb A^1} (\varphi_*\pr^*_C\beta) \\
    &= 0^!_{(E_1\oplus E_2) \times \mathbb A^1} (\varphi_*(\beta \times [\mathbb A^1])) \in A^T_*(\mathcal Z\times\mathbb A^1).
  \end{align*}
  For each \(t\), the inclusion \(r_t:\{t\}\hookrightarrow \mathbb A^1\) yields Kresch's refined Gysin map \(r^!_t\). It acts a map \(r^!_t:A^T_*(Y)\to A^T_*(Y_t)\) for any \(Y\) over \(\mathbb A^1\), it commutes projective pushforward and flat pullback and \(r^!_t\circ \pr_C = \id\) on \(A^T_*(C)\). Therefore,
  \begin{align*}
    r^!_t\Theta(\beta) &= r^!_t0^!_{(E_1\oplus E_2) \times \mathbb A^1}(\varphi_*(\beta \times [\mathbb A^1])) \\
    &= 0^!_{E_1\oplus E_2}(r^!_t\varphi_{*}(\beta \times [\mathbb A^1])) \\
    &= 0^!_{E_1\oplus E_2}(\varphi_{t*}(\beta)).
  \end{align*}
  Since \(r^!_t\Theta(\beta) = (\pr_C)^{-1}\circ \Theta(\beta)\) is independent of \(t\), so is \(0^!_{E_1\oplus E_2}(\varphi_{t*}(\beta))\), and hence
  \begin{align*}
    0^!_{E_1\oplus E_2} (i_1\beta) = 0^!_{E_1\oplus E_2}(\varphi_{1*}(\beta)) = 0^!_{E_1\oplus E_2}(\varphi_{0*}(\beta)) = 0^!_{E_1\oplus E_2}(i_0\beta).
  \end{align*}

  All that is left is to transport the inverse Euler class through the various morphisms. Let \(f:E_1\oplus E_2\to \mathcal Z\) be the bundle map and factor it as \(f = g\circ h\) where \(h:E_1\oplus E_2\to E_1\) is the first projection and \(g:E_1\to \mathcal Z\) is the bundle map of \(E_1\). This means
  \begin{align*}
    0^!_{E_1\oplus E_2} = (f^*)^{-1} = (g^{*})^{-1}\circ (h^*)^{-1} = 0^!_{E_1} \circ (h^*)^{-1}.
  \end{align*}
  Now remember that we defined \(\iota:E_1\to E_1\oplus E_2\), so the self-intersection formula for a class \(\zeta\) in \(A_*^T(E_1)\) reads
  \begin{align*}
    (h^*)^{-1}(\iota_* \zeta) = e_T(g^*E_2) \cap \zeta,
  \end{align*}
  which we apply below to \(\zeta = j_{1*}\beta\).

  Applying the Whitney sum formula and the projection formula (\cite[\S 3.6]{kresch_CycleGroupsArtin1999}) then yields
  \begin{align*}
    S_{(i_1, E_1\oplus E_2)}(\beta) &= e^{-1}_T(E_1\oplus E_2) \cap 0^!_{E_1\oplus E_2}(i_{1*}\beta) \\
    &= e^{-1}_T(E_1)\cdot e_T^{-1}(E_2) \cap 0^!_{E_1\oplus E_2}(i_{0*}\beta) \\
    &= e^{-1}_T(E_1)\cdot e_T^{-1}(E_2) \cap 0^!_{E_1}\big((h^*)^{-1}(\iota_*j_{1*}\beta)\big) \\
    &= e^{-1}_T(E_1)\cdot e_T^{-1}(E_2) \cap 0^!_{E_1}\big(e_T(g^*E_2) \cap j_{1*}\beta\big) \\
    &= e^{-1}_T(E_1)\cdot e_T^{-1}(E_2)\cdot e_T(E_2) \cap 0^!_{E_1}\big(j_{1*}\beta\big) \\
    &= e_T^{-1}(E_1) \cap 0^!_{E_1}\big(j_{1*}\beta \big) \\
    &= S_{(j_1, E_1)}(\beta).
  \end{align*}
  Swapping the roles of \(E_1\) and \(E_2\) and using the interpolation
  \begin{align*}
    \varphi'(c,t) = (t·j1(c), j2(c), t)
  \end{align*}
  gives us \(S_{i_1, E_1\oplus E_2}(\beta) = S_{j_2, E_2}(\beta)\), and we are done.
\end{proof}


This lemma fails for cone stacks, which is why it isn't a proof of Lemma~\ref{lem:independence-of-Res-map-from-choice} on its own. The following example demonstrates what goes wrong:
\begin{example}\label{rl-0012}
  Suppose \((j_i:\mathfrak C\to \mathfrak E_i)\) is an embedding of the intrinsic normal cone stack \(\mathfrak C\) for \(\eta:Z\to X\) into a vector bundle stack \(\mathfrak E_i\) where \(i = 1,2\). Then it is not necessarily the case that the diagonal map \(j=(j_1,j_2):\mathfrak C\to \mathfrak E_1\oplus \mathfrak E_2\) is a closed embedding.

  To see this, let \(Z = X = \Spec k\), \(V = \mathbb A^1\) be the rank 1 vector bundle over \(Z\) and \(\mathfrak C = \mathfrak E_1 = \mathfrak E_2 = [0 / V] = BV\). Note that \(\mathfrak C\) is the intrinsic normal cone of the morphism \(\Spec k \to [\Spec k/V]\). Take \(j_i:\mathfrak C\to \mathfrak E_i\) to be the identity for both \(i\). The diagonal is then
  \[j = (j_1,j_2): BV \to BV \oplus BV = B(V\oplus V)\]
  induced by the diagonal map \(V\to V\oplus V\). The fiber of a point \(\Spec k\) over this diagonal is
  \[\Spec k\times_{B(V\oplus V)}BV = V,\]
  so the map \(j\) has positive dimensional fibers.

  On the level of complexes, this can be seen as a failure of \(h^0\) to be an isomorphism (as per the criterion in \cite[Proposition 2.6]{behrend-fantechi_IntrinsicNormalCone1997}). The two-term truncation of the cotangent complex of \(f:\Spec k \to BV\) is \(L_f = [0 \to k]\) (in degree \(-1\) and \(0\)) and the complex corresponding to \(B(V\oplus V)\) is \([0 \to k^2]\) (dualize and quotient to get \(B(V\oplus V)\)). The candidate obstruction theory is then determined by the morphism in degree \(0\), which is \((a,b) \mapsto a+b\), and this is exactly the map \(h^0(\phi)\) as well. This is surjective, but not an isomorphism, hence \(\phi^\vee:h^1/h^0(L_f^\vee)\to h^1/h^0\big((E_1^\bullet\oplus E_2^\bullet)^\vee\big)\) fails to be a closed embedding.
\end{example}


We fix this issue by moving from stacky cones and vector bundles to their atlases. Given two choices of stacky embedding data \((j_i, \mathfrak E_i)\) with \(i=1,2\) for the normal cone \(\mathfrak C_\eta\) where \(\mathfrak E_i = [E^1_i/E^0_i]\), set
\begin{align*}
  C_{E^1_i} = \mathfrak C_i\times_{\mathfrak E_i}E^1_i
\end{align*}
and form \(j'_i:C_{E_i^1}\hookrightarrow E_i^1\) by pulling \(j_i\) back along the atlas \(p_i:E^1_i\to \mathfrak E_i\). This \(C_{E^1_i}\) is an honest \(T\)-cone over \(\mathcal Z\); smooth locally on \(\mathcal Z\) the cone stack \(\mathfrak C\) is a quotient of a cone by a bundle action, and the atlas pullback recovers that cone. We then show that
\begin{align*}
  S_{(j_1, \mathfrak E_1)} = S_{(j_2, \mathfrak E_2)}
\end{align*}
by comparing \(S_{C_{E^1_1}}\) and \(S_{C_{E^1_2}}\). We need one more preliminary result, after which \ref{lem:independence-of-Res-map-from-choice} is immediate.

\begin{lemma}\label{rl-0013}\label{lem:torsor-lemma}
  Suppose
  \begin{align*}
    0 \to F\to D\xrightarrow {p} \mathfrak C\to 0
  \end{align*}
  is an exact sequence of moving cone stacks over \(\mathcal Z\) where \(F\) is a finite rank moving vector bundle, \(D\) is a moving cone and \(\mathfrak C\) is a moving cone stack. Then
  \begin{align*}
    S_C(\beta) = e_T(F)\cap S_D(p^*\beta).
  \end{align*}
\end{lemma}
\begin{proof}
  If \(j:D\hookrightarrow G\) is an embedding of \(D\) into a finite rank moving vector bundle stack over \(\mathcal Z\), then we get a diagram
  \begin{center}
    \begin{tikzcd}[column sep={1.5cm,between origins},
      row sep={1.5cm,between origins}]
      0 \arrow[r] & F \arrow[r] & G \arrow[r, "p'"] & \mathfrak E \arrow[r] & 0 \\
      0 \arrow[r] & F \arrow[u, hook, "\id"] \arrow[r] & D \arrow[u, hook, "j"] \arrow[r,"p"] & \mathfrak C \arrow[u, hook, "j'"] \arrow[r] & 0
    \end{tikzcd}.
  \end{center}
  The properties of Gysin morphisms and flat pullback mean that for any \(\beta\in A^T_*(C)\) we have \(jp^*(\beta) = p'^*j'_*(\beta)\) and \(0^!_G\circ p'^* = 0^!_E\), so
  \begin{align*}
    S_C(\beta) &= e_T(\mathfrak E)^{-1} \cap 0^!_{E}j_*(\beta) \\
           &= e_T(F)\cdot e_T(F)^{-1}\cdot e_T(\mathfrak E)^{-1}~\cap~ 0^!_Gp'^*j_*(\beta) \\
           &= e_T(F)\cdot e_T(G)^{-1} ~\cap~ 0^!_Gj'_*(p^*\beta) \\
    &= e_T(F)\cap S_D(p^*\beta).
  \end{align*}
\end{proof}


The independence lemma is now straightforward.

\begin{proof}[Proof of Lemma~\ref{lem:independence-of-Res-map-from-choice}]
  Let \((j_1, \mathfrak E_1)\) and \((j_2, \mathfrak E_2)\) be two choices of embedding data for the intrinsic normal cone \(\mathfrak C_\eta\) of \(\eta:\mathcal Z\to \mathcal X\). We wish to show
  \[S_{(j_1, \mathfrak E_1)}(\beta) = S_{(j_2, \mathfrak E_2)}(\beta)\] for all \(\beta \in A^T_*(\mathfrak C_\eta)\).
  Let \(\mathfrak E_i = [E^1_i/E^0_i]\) be presentations of these vector bundle stacks as global quotients for \(i = 1,2\). Form the following diagram:
  \begin{center}
    \begin{tikzcd}[column sep={1.5cm,between origins},
      row sep={1.5cm,between origins}]
      & & C \arrow[dl,swap,two heads,"r_1"] \arrow[dr,two heads,"r_2"] \arrow[dd, phantom, "\square" {rotate=45}]& & \\
      & C_{E^1_1} \arrow[dd,phantom,"\square" {rotate=45}] \arrow[dl,swap,hook',"j_1'"] \arrow[dr,swap,two heads, "p'_1"] & & C_{E^1_2} \arrow[dd,phantom,"\square" {rotate=45}] \arrow[dl,two heads,"p'_2"] \arrow[dr,hook,"j_2'"] & \\
      E^1_1 \arrow[dr,swap,two heads,"p_1"] & & \mathfrak C_\eta \arrow[dl,hook',"j_1"] \arrow[dr,swap,hook,"j_2"] & & E_2^1 \arrow[dl,two heads,"p_2"] \\
      & \mathfrak E_1 & & \mathfrak E_2 & 
    \end{tikzcd}
  \end{center}
  For each \(i=1,2\) the projection \(p'_i:C_{E^1_i}\to \mathfrak C_\eta\) is an \(E^0_i\)-torsor over \(\mathfrak C_\eta\). The fiber product \(C = C_{E^1_1}\times_{\mathfrak C_\eta}C_{E^1_2}\) is an \(E^0_2\)-torsor over \(C_{E^1_1}\) via \(r_1\) and an \(E^0_1\)-torsor over \(C_{E^1_2}\) via \(r_2\), since torsors are compatible with base change.

  By Lemma \ref{lem:torsor-lemma},
  \begin{equation}\label{eqn:independence-lemma-eqn1}
    S_{(j_i, \mathfrak E_i)}(\beta) = e_T(E^0_i) \cap S_{C_{E^1_i}}(p'^*_i\beta)
  \end{equation}
  for \(i = 1,2\). Then by repeated application of this formula and Lemma~\ref{lem:torsor-lemma}~we have
  \begin{align*}
    S_{(j_1, \mathfrak E_1)}(\beta) &= e_T(E^0_1) \cap S_{C_{E^1_1}}(p'^*_1(\beta)) \tag{Equation \eqref{eqn:independence-lemma-eqn1}}\\
    &= e_T(E^0_1) \cap e_T(E^0_2) \cap S_{C}(r_1^*p'^*_1(\beta)) \tag{Lemma~\ref{lem:torsor-lemma}}\\
    &= e_T(E^0_2) \cap e_T(E^0_1) \cap S_{C}(r_2^*p'^*_2(\beta)) \tag{top square commutativity}\\
    &= e_T(E^0_2) \cap S_{C_{E^1_2}}(p'^*_2(\beta)) \tag{Lemma~\ref{lem:torsor-lemma}}\\
    &= S_{(j_2, \mathfrak E_2)}(\beta), \tag{Equation \eqref{eqn:independence-lemma-eqn1}}
  \end{align*}
  and we're done.
\end{proof}

\begin{definition}\label{rl-0014}\label{defn:residue-map}
  Let \(\mathcal X\) be an algebraic stack of finite type over \(\Spec k\) with a \(T\)-action and suppose that \(\mathcal X^T\) is algebraic. Let \(\mathcal Z\subseteq \mathcal X^T\) be an open and closed substack which is finite type over \(k\) and has the resolution property. We then define the equivariant residue \(\Res_{\mathcal Z/\mathcal X}:A^T_*(\mathcal X)\to A^T_*(\mathcal Z)_{\loc}\) by
  \begin{align*}
    \Res_{\mathcal Z/\mathcal X}(\alpha) = e^{-1}_T(\mathfrak E) \cap \eta^!(\alpha),
  \end{align*}
  where \(\mathfrak E\) is any moving equivariant vector bundle stack over \(\mathcal Z\) admitting an embedding \(j:\mathfrak C_{\mathcal X^T/\mathcal X}\hookrightarrow \mathfrak E\).
\end{definition}

